\documentclass[10pt,a4paper]{amsart}
\usepackage[margin=19mm]{geometry}
\usepackage{amsmath,amssymb,mathtools}
\usepackage[scr=rsfs,cal=euler]{mathalfa}
\usepackage{xfrac}
\usepackage[unicode,hidelinks,bookmarks=false]{hyperref}
\newcommand{\ie}{i.e.\ }
\newcommand{\cf}{cf.\ }
\newcommand{\resp}{respectively\ }
\newcommand{\Subsection}{\subsection}
\providecommand{\nobreakdash}{\nobreak\hbox{-}}
\setlength{\emergencystretch}{2em}
\multlinegap0pt
\usepackage{mathtools}
\usepackage{float}
\newtheorem{theorem}{Theorem}
\newtheorem{lemma}{Lemma}
\newtheorem{corollary}{Corollary}
\renewcommand\P{\mathbb{P}}
\newcommand\Z{\mathbb{Z}}
\newcommand\be{\mathbf{e}}
\DeclareMathOperator{\card}{card}
\newcommand\eps{\varepsilon}
\title{Two-dimensional Rademacher walk}
\author{Satyaki Bhattacharya, Stanislav Volkov}
\date{Source: arXiv:2506.16259v2 (CC BY 4.0)}
\begin{document}
\maketitle

\noindent\emph{Source and licence.} Two-dimensional Rademacher walk. Version arXiv:2506.16259v2 (CC BY 4.0), distributed by arXiv under the Creative Commons Attribution 4.0 International licence, http://creativecommons.org/licenses/by/4.0/, as recorded in the arXiv licence metadata for this exact version and shown on the abstract page https://arxiv.org/abs/2506.16259v2. Full licence notice: \texttt{licenses/arxiv-2506.16259-CC-BY-4.0.txt}.\par\smallskip

\noindent\textit{Editorial scope.} Counted unit: the article's theorem t1 (source0.tex lines 73--75). The article's complete original proof of it is delivered with it (source0.tex lines 130--189, one \texttt{proof} environment of 60 lines, which the article places away from the statement). No mathematics is written, repaired or re-derived: the statement and the proof are the article's own lines, in the article's own notation, and no retained line is substituted or re-flowed.  Retained uncounted as the source-local context the proof needs: lines 78--84; lines 106--112.  Nothing was omitted from any retained span, so the package carries no omission record.  No retained line contains a citation, so the wrapper carries no bibliography and no bibliography entry.  No retained line contains a figure environment, an image-inclusion command or drawing code, so no image is delivered and the package declares no asset dependency.  Editorial formatting only, in addition: an AMS \texttt{amsart} preamble that restates the article's own declarations and macros used by the retained lines, namely the environments \texttt{theorem}, \texttt{lemma}, \texttt{corollary} on separate counters, together with the standard packages those lines need.  The retained environments print in this document as Theorem 1, Lemma 1, Corollary 1 in order of appearance, whereas the article prints the same items with the numbering of its own full document.  The complete native source of the article as distributed in the arXiv e-print for this exact version is bundled unchanged as \texttt{sources/180143/source0.tex}.
\begin{lemma}\label{modlemma}
Let $d_1,d_2,d_3\dots$ be a sequence of distinct positive integers, and let $T_k=d_1\eps_1+\dots+d_k\eps_k$ where  $\eps_i$ are i.i.d.\ Rademacher. Then for every integer $\displaystyle  m\ge \max_{1\le i\le k}d_i$ and every  $M\in\Z$ we have
$$
\P(T_k-M\equiv 0\pmod m)\le \frac{C \log k}{k}    
$$
for  some universal constant $C>0$.
\end{lemma}

\begin{corollary}\label{modcor}
Let $d_1,d_2,d_3\dots,d_n$ be a sequence of positive integers, of which at least $k$ are distinct, say $d_{\ell_1},d_{\ell_2},\dots,d_{\ell_k}$, and let $T_n=\sum_{i=1}^n d_i\eps_i$ where  $\eps_i$ are i.i.d.\ Rademacher. Then for every integer $\displaystyle  m\ge \max_{1\le i\le k}d_{\ell_i}$ we have
$$
\P(T_n\equiv 0\pmod m)\le \frac{C \log k}{k}    
$$
where $C$ is the constant from the statement of Lemma~\ref{modlemma}.
\end{corollary}

\begin{theorem}\label{t1}
Suppose that $a_n$ is a non-decreasing sequence of integers converging to infinity. Then $S_n$ is transient.
\end{theorem}

\begin{proof}[Proof of Theorem~\ref{t1}]
Let $b_1=a_1\ge 1$, $b_2=\min\{a_n: a_n>b_1\}\ge a_1+1$,  $b_3=\min\{a_n: a_n>b_2\}\ge a_2+1$, etc. Denote by $m_i=\card\{n:\ a_n=b_i\}$, $i=1,2,\dots$. Then the sequence $\{a_n\}$ can be represented as
$$
\underbrace{b_1, b_1, \dots b_1,}_{m_1}
\underbrace{b_2, b_2, \dots b_2,}_{m_2} 
\underbrace{b_3, b_3, \dots b_3,}_{m_3} 
\dots
\underbrace{b_j, b_j, \dots b_j,}_{m_j}b_{j+1},\dots .
$$
Let  $k_{j+1}=m_1+m_2+\dots+m_{j}+1=\inf\{n:\ a_n=b_{j+1}\}$ be the first index of the sequence $\{a_n\}$ equal to $b_{j+1}$.

Let $\kappa_i=1$ if $\xi_i\in\{\be_1,-\be_1\}$ and $=0$ otherwise. Let $\eps_i=\pm 1$ equal to the non-zero coordinate of~$\xi_i$, Then $\{\kappa_i\}$ are i.i.d.\ Bernoulli and $\{\eps_i\}$ are i.i.d.\ Rademacher; moreover, these two sequences are independent. The values of $i\in\{1,2,\dots,n\}$ for which $\kappa_i=1$ ($=0$ resp.) correspond to horizontal (vertical, resp.) steps. Formally, let
$$
H_n=\{1\le i<n:\ \kappa_i=1\},\quad
V_n=\{1\le i< n:\ \kappa_i=0\}
$$
be the indices corresponding to horizontal (resp.\ vertical) steps. Note that $Z_i=X_i \be_1+Y_i\be_2$ where 
$$
X_i=\begin{cases}
   a_i\eps_i,&\text{if } \kappa_i=1;\\
   0, &\text{if } \kappa_i=0
\end{cases}
\qquad
Y_i=\begin{cases}
   0,&\text{if } \kappa_i=1;\\
   a_i\eps_i,&\text{if } \kappa_i=0,
\end{cases}
$$
and denote 
$$
S_n=S^X_n\be_1+S^Y_n\be_2\quad \text{ where }
S_n^X=X_1+\dots+X_n,\quad S_n^Y=Y_1+\dots+Y_n.
$$
Fix a large integer $j$ and let
$$
A_j=\left\{\card\{a_i,i\in H_{k_j}\}<\frac j3\text{ \ or \ }\card\{a_i,i\in V_{k_j}\}<\frac j3\right\}.
$$
For each $1\le i< k_j$, $\kappa_i$ equals $0$ or $1$ with equal probability $\frac12$. Hence, by Hoeffding's inequality, there exists a constant $c_0$ such that  
$$
\P(A_j)\le 2e^{-c_0 j},
$$ 
i.e., there are at least $j/3$ horizontal and $j/3$ vertical steps with probability tending to one. As a result, on $A_j^c$, we can write
$$
S^X_{k_j-1}=\sum_{i=1}^{h_j} b'_i \eps'_i, \quad S^Y_{k_j-1}=\sum_{i=1}^{v_j} b''_i \eps''_i
$$
where $h_j=\card\{H_{k_j}\}$, $v_j=\card\{V_{k_j}\}$, the positive integer sequences $\{b_i'\}_{i=1}^{h_j}$ and $\{b_i''\}_{i=1}^{v_j}$ are non-decreasing and each has at least $j/3$ {\em distinct} elements, while $\eps_i',\eps_i''$ are i.i.d.\ Rademacher.

Now, by Corollary~\ref{modcor},
\begin{align*}
\P(S^X_{k_j-1}\equiv 0\pmod {b_j} \mid A_j^c)\le \frac{C\log (j/3)}{j/3},\qquad
\P(S^Y_{k_j-1}\equiv 0\pmod {b_j} \mid A_j^c)\le \frac{C\log (j/3)}{j/3}
\end{align*}
and, moreover, conditioned on the sequence $\kappa_i$ and $A_j^c$, the events above are independent. Since for the walk to hit $\mathbf{0}$ at time $n$, $k_{j}\le n\le k_{j+1}-1$, we need that both its coordinates are divisible by~$b_j$ at time~$k_j-1$, we have
\begin{align}\label{eqBC}
\P(S_n=\mathbf{0}\text{ for some }n\in [k_j,k_{j+1}))
\le \P(A_j)+\left(\frac{C\log (j/3)}{j/3}\right)^2
=O\left(\frac{\log^2 j}{j^2}\right).
\end{align}
Since the quantity in~\eqref{eqBC} is summable in $j$, by the Borel-Cantelli lemma we conclude that the event $\{S_n=\mathbf{0}\}$ occurs finitely often. By the same token, we get that $\{S_n=\mathbf{w}\}$ occurs finitely often for each ${\bf w}\in \Z^2$. Hence, the transience follows.
\end{proof}

\end{document}
